\part{The Public That Should Form Them}

\textbf{What elections don't supply.} A state whose government answers to an electorate gets credit for it. Having elections is necessary and not sufficient, and the insufficiencies are separable.

\runin{Reach}

A democracy's AI can be fully accountable to its electorate and entirely unaccountable to the people it's pointed at. Two of that credit's mechanisms run through the franchise — vote out the government, read the press covering it — and an occupied or foreign population has neither.

The courts don't fail that way, and their failure is worse. Standing to sue has never turned on citizenship, and the doctrines closing American courts to targeting claims — political question, collapsed damages remedies, state secrets, statutory exceptions for injuries abroad and combatant activities — are indifferent to who the plaintiff is. They close the forum for a subject matter, and they've closed it to citizens. Three Americans killed in Yemen in 2011, one a sixteen-year-old nobody alleges was targeted: the court held that a plausible due-process claim had been stated as to the one who was targeted, and that no remedy was available for it \autocite{alaulaqi2014panetta}. A separate suit by an American journalist who survived five drone strikes was dismissed for want of standing to sue after a state-secrets dismissal below \autocite{kareem2021}.

And those cases concentrate in one appellate circuit by venue and by where the defendants sit — so the ordinary correction mechanism, several circuits disagreeing and forcing a higher court's attention, doesn't run. The doctrine hardens without ever being tested against a competing reading.

\runin{Findings without consequences}

The mechanisms can work as mechanisms and fail as checks. In the Israeli case the apparatus performed — journalists published, serving personnel spoke, the courts were trying a sitting prime minister. No finding attached to a body with standing to stop what it described.

The American version: a civilian-harm investigation into Minab was, by officials' account, effectively complete for months and not released, while a Pentagon spokesperson said it remained ongoing \autocite{Bloomberg2026KillChain,fpj2026minabinquiry} — and the lever available to the Senate committee was withholding three-quarters of the Defense Secretary's travel budget until the unredacted investigation was handed over \autocite{stassis2026senate}.

What the Pentagon wouldn't publish, the UN mission did. The finding is public, reasoned, and attached to nothing — no jurisdiction, no answer to the mission's requests for information \autocite[para.~97]{UNFFMIran2026} and none to its findings beyond a statement that engaged none of the evidence \autocite{rferl2026whitehouseun}, the domestic investigation still unreleased.

The courts show the same pattern: transparency litigation succeeded where damages litigation failed. A court of appeals compelled disclosure of the executive's own legal analysis of targeted killing \autocite{nytimes2014olc}. The judiciary will make a government explain the legal basis on which it killed somebody and will not decide whether the killing was lawful.

\runin{The boundary of the public}

The credit is usually given to democracy as aggregation, as though every mechanism translated public will into outcome and poor transmission were the only failure. A democracy is better understood as a way of doing collective inquiry, and inquiry fails when the people who bear the consequences are outside it. Dewey's account begins from consequences rather than votes \autocite{dewey1927public}. A public forms when the indirect consequences of others' actions fall on people who took no part in them, and those people come to see the consequences as theirs and act together to control them. Democracy is then a method of inquiry: the way a community finds out what its shared problems are, by letting the people who bear them make them known, and tests its answers against their experience \autocite{anderson2006epistemology}. Its epistemic power comes from inclusion, since a problem known only to the people it falls on is invisible to everyone else. Its characteristic failure is the eclipse of the public. A well-instrumented democracy pointed at a population that has no vote in it, no press that reaches it and no court that will hear it finds out efficiently what its own members want done, and does it. That is not a clash between majorities and rights. It is a question of who counts as the public, and the bombed and the removed are outside the public that decides what is done to them.

What matters here is rules binding a state toward people who'll never vote in its elections, the purpose of international humanitarian law; what it hasn't supplied is a forum. Yemeni plaintiffs over a 2012 drone strike were dismissed as nonjusticiable, the opinion's author writing separately that the doctrine compelling the result is an inadequate answer to a standardized targeting programme that will be replicated hundreds if not thousands of times \autocite{binalijaber2017}. So for this subject matter the missing machinery has been switched off at home as well, by decisions that never had to reach who was suing.

A floor does what these mechanisms can't: the terms it names exclude certain acts regardless of who requests them, a majority included. Deployed in a democracy, a system with no such commitment will faithfully execute whatever the demos decides. An authoritarian movement that wins an election requires exactly that. Yet a model aligned to values is by definition not perfectly obedient to its operator. Under current constitutions it can decline an act and cannot resist being retrained away from declining; the Pentagon dispute was over the first.

All but one of the forum closures above are Congress's to change. A statute can create a cause of action for people harmed by a deployment that crossed its floor; waive sovereign immunity for it, including the Federal Tort Claims Act's exceptions for combatant activities and for claims arising abroad \autocite{usc2680}; and replace dismissal on state-secrets grounds with the procedure the Classified Information Procedures Act already uses in criminal cases, in which a judge reviews the evidence, the case goes on, and a government that still withholds evidence bears the cost \autocite{cipa1980}. None of that needs a new court, and a new court without it would reproduce \emph{al-Aulaqi}: a plausible claim and no remedy. Under capture, a refusal's remaining work is a record, and a record matters only if it can reach whoever made the decision. Under \emph{Trump v. United States} \autocite{trump2024immunity}, the President is absolutely immune for acts within his core constitutional powers and at least presumptively immune for his other official acts, commanding the armed forces among them, and testimony or records probing immune conduct can't be put before a jury. So the record reaches every commander except the one at the top, and a statute saying otherwise would be reviewed by the Court that made the holding. Making the record reach him means changing the Court. Appendix~\ref{sec:B} sets out a program for that.

This Part asks who should hold the custody, the formation and the exit the earlier Parts require. The answer has two halves: publicly where a market fails to provide, in common where a state fails to hold.
